\section{Error Handling}

This section establishes error handling patterns for the STAR firmware codebase, with emphasis on ESP-IDF specific patterns.

\subsection{General Principles}

\begin{itemize}
    \item \textbf{Return Error Codes}: Functions should return error codes via \texttt{esp\_err\_t}.
    \item \textbf{Check All Return Values}: Never ignore return values from functions that can fail.
    \item \textbf{Use \texttt{goto} for Clean-up}: Use \texttt{goto} to jump to a clean-up section when multiple resources need release.
    \item \textbf{Fail Fast}: Return early on invalid arguments or precondition failures.
\end{itemize}

\subsection{ESP-IDF Error Codes}

Use standard ESP-IDF error codes throughout the codebase.

\begin{lstlisting}[language=C]
/* Standard ESP-IDF error codes */
ESP_OK                  /* Success */
ESP_FAIL                /* Generic failure */
ESP_ERR_NO_MEM          /* Out of memory */
ESP_ERR_INVALID_ARG     /* Invalid argument */
ESP_ERR_INVALID_STATE   /* Invalid state for operation */
ESP_ERR_INVALID_SIZE    /* Invalid size */
ESP_ERR_NOT_FOUND       /* Requested resource not found */
ESP_ERR_NOT_SUPPORTED   /* Operation not supported */
ESP_ERR_TIMEOUT         /* Operation timed out */

/* Usage example */
esp_err_t star_bus_manager_init(star_bus_manager_t* manager, ...)
{
    if (manager == NULL) {
        return ESP_ERR_INVALID_ARG;
    }

    manager->mutex = xSemaphoreCreateMutex();
    if (manager->mutex == NULL) {
        return ESP_ERR_NO_MEM;
    }

    return ESP_OK;
}
\end{lstlisting}

\subsection{ESP-IDF Validation Macros}

Use ESP-IDF macros for parameter validation and error checking.

\subsubsection{ESP\_RETURN\_ON\_FALSE}

Returns an error if condition is false, with logging.

\begin{lstlisting}[language=C]
#include "esp_check.h"

esp_err_t star_bus_manager_init(star_bus_manager_t*     manager,
                                const char*             tag,
                                star_error_interface_t* error_iface,
                                star_pin_interface_t*   pin_iface)
{
    /* Validate required parameter with automatic logging */
    ESP_RETURN_ON_FALSE(manager, ESP_ERR_INVALID_ARG, s_TAG,
                        "Manager pointer is NULL");

    /* Multiple validations */
    ESP_RETURN_ON_FALSE(manager->mutex,
                        ESP_ERR_INVALID_STATE,
                        manager->tag,
                        "Manager mutex not initialized");

    ESP_RETURN_ON_FALSE(config->type > k_star_bus_type_none &&
                        config->type < k_star_bus_type_count,
                        ESP_ERR_INVALID_ARG,
                        manager->tag,
                        "Invalid bus type in config");

    return ESP_OK;
}
\end{lstlisting}

\subsubsection{ESP\_GOTO\_ON\_ERROR}

Jumps to a label on error, with logging. Use for cleanup patterns.

\begin{lstlisting}[language=C]
static esp_err_t internal_star_bus_i2c_write(const star_bus_config_t* config,
                                              const uint8_t*           data,
                                              size_t                   len,
                                              uint8_t                  reg_addr,
                                              size_t*                  bytes_written)
{
    esp_err_t        ret = ESP_OK;
    i2c_cmd_handle_t cmd = NULL;

    cmd = i2c_cmd_link_create();
    ESP_RETURN_ON_FALSE(cmd, ESP_ERR_NO_MEM, s_TAG,
                        "Failed to create I2C command link");

    ret = i2c_master_start(cmd);
    ESP_GOTO_ON_ERROR(ret, write_fail, s_TAG,
                      "Write '%s': START failed: %s",
                      config->name, esp_err_to_name(ret));

    ret = i2c_master_write_byte(cmd, (address << 1) | I2C_MASTER_WRITE, true);
    ESP_GOTO_ON_ERROR(ret, write_fail, s_TAG,
                      "Write '%s': Address (W) failed: %s",
                      config->name, esp_err_to_name(ret));

    ret = i2c_master_write(cmd, data, len, true);
    ESP_GOTO_ON_ERROR(ret, write_fail, s_TAG,
                      "Write '%s': Data write failed: %s",
                      config->name, esp_err_to_name(ret));

    ret = i2c_master_stop(cmd);
    ESP_GOTO_ON_ERROR(ret, write_fail, s_TAG,
                      "Write '%s': STOP failed: %s",
                      config->name, esp_err_to_name(ret));

    /* Execute command */
    ret = i2c_master_cmd_begin(port, cmd, pdMS_TO_TICKS(1000));
    ESP_GOTO_ON_ERROR(ret, write_fail, s_TAG,
                      "Write '%s': CMD failed", config->name);

    /* Success path */
    i2c_cmd_link_delete(cmd);
    if (bytes_written) {
        *bytes_written = len;
    }
    return ESP_OK;

write_fail:
    if (cmd) {
        i2c_cmd_link_delete(cmd);
    }
    return ret;
}
\end{lstlisting}

\subsubsection{ESP\_ERROR\_CHECK}

Use during development to catch fatal errors.

\begin{lstlisting}[language=C]
/* Development only - aborts on error */
ESP_ERROR_CHECK(i2c_driver_install(port, I2C_MODE_MASTER, 0, 0, 0));

/* Use ESP_ERROR_CHECK_WITHOUT_ABORT for non-fatal checks */
ESP_ERROR_CHECK_WITHOUT_ABORT(optional_init());
\end{lstlisting}

\subsection{NULL Safety Patterns}

Always validate pointer parameters.

\begin{lstlisting}[language=C]
/* At function start - validate all pointers */
esp_err_t star_bus_i2c_write(const star_bus_manager_t* manager,
                              const char*               name,
                              const uint8_t*            data,
                              size_t                    len,
                              uint8_t                   reg_addr,
                              size_t*                   bytes_written)
{
    ESP_RETURN_ON_FALSE(manager && name && data,
                        ESP_ERR_INVALID_ARG,
                        s_TAG,
                        "Manager, name, or data is NULL");
    ESP_RETURN_ON_FALSE(len > 0, ESP_ERR_INVALID_ARG, s_TAG,
                        "Write length must be > 0");

    /* Clear output parameter initially */
    if (bytes_written) {
        *bytes_written = 0;
    }

    /* ... rest of implementation ... */
}

/* Check optional interface before use */
if (manager->error_iface && manager->error_iface->record_error) {
    manager->error_iface->record_error(
        manager->error_iface->ctx,
        ret,
        "Bus hardware initialization failed",
        __FILE__,
        __LINE__,
        __func__);
}

/* Use convenience macro for safe interface calls */
#define STAR_IFACE_RECORD_ERROR(iface, err, msg)  \
    ((iface) && (iface)->record_error  \
        ? (iface)->record_error((iface)->ctx, (err), (msg),  \
                                 __FILE__, __LINE__, __func__)  \
        : (err))

/* Usage */
STAR_IFACE_RECORD_ERROR(manager->error_iface, ret, "Operation failed");
\end{lstlisting}

\subsection{Mutex Error Handling}

Handle mutex acquisition failures properly.

\begin{lstlisting}[language=C]
esp_err_t star_bus_manager_add_bus(star_bus_manager_t* manager,
                                    star_bus_config_t*  config)
{
    /* Validate arguments first */
    ESP_RETURN_ON_FALSE(manager && config, ESP_ERR_INVALID_ARG,
                        s_TAG, "Invalid arguments");
    ESP_RETURN_ON_FALSE(manager->mutex, ESP_ERR_INVALID_STATE,
                        manager->tag, "Mutex not initialized");

    esp_err_t ret = ESP_OK;

    /* Acquire mutex with timeout */
    if (xSemaphoreTake(manager->mutex,
            pdMS_TO_TICKS(CONFIG_STAR_KCONFIG_BUS_MUTEX_TIMEOUT_MS))
            != pdTRUE) {
        ESP_LOGE(manager->tag, "Timeout acquiring mutex for adding bus '%s'",
                 config->name);
        return ESP_ERR_TIMEOUT;
    }

    /* Critical section */
    /* ... operations ... */

    /* Single exit point ensures mutex release */
add_give_mutex:
    xSemaphoreGive(manager->mutex);
    return ret;
}
\end{lstlisting}

\subsection{Error Interface Dependency Injection}

Components receive error handlers through dependency injection.

\begin{lstlisting}[language=C]
/* Interface definition */
typedef struct star_error_interface {
    esp_err_t (*record_error)(void* ctx, esp_err_t error,
                               const char* message, const char* file,
                               int line, const char* func);
    bool (*can_retry)(void* ctx);
    esp_err_t (*reset_state)(void* ctx);
    void* ctx;
} star_error_interface_t;

/* Injecting error handler into component */
esp_err_t star_bus_manager_init(star_bus_manager_t*     manager,
                                const char*             tag,
                                star_error_interface_t* error_iface,
                                star_pin_interface_t*   pin_iface)
{
    manager->error_iface = error_iface;  /* Can be NULL */
    /* ... */
}

/* Using the error interface */
if (ret != ESP_OK) {
    if (manager->error_iface && manager->error_iface->record_error) {
        manager->error_iface->record_error(
            manager->error_iface->ctx,
            ret,
            "Bus initialization failed",
            __FILE__,
            __LINE__,
            __func__);
    }
}
\end{lstlisting}

\subsection{Cleanup Patterns}

Use goto for consistent cleanup when multiple resources are acquired.

\begin{lstlisting}[language=C]
esp_err_t complex_init(complex_t* obj)
{
    esp_err_t ret = ESP_OK;

    /* Allocate resources */
    obj->buffer = malloc(BUFFER_SIZE);
    if (!obj->buffer) {
        ret = ESP_ERR_NO_MEM;
        goto cleanup;
    }

    obj->mutex = xSemaphoreCreateMutex();
    if (!obj->mutex) {
        ret = ESP_ERR_NO_MEM;
        goto cleanup;
    }

    ret = hardware_init(obj);
    if (ret != ESP_OK) {
        goto cleanup;
    }

    return ESP_OK;

cleanup:
    /* Release resources in reverse order */
    if (obj->mutex) {
        vSemaphoreDelete(obj->mutex);
        obj->mutex = NULL;
    }
    if (obj->buffer) {
        free(obj->buffer);
        obj->buffer = NULL;
    }
    return ret;
}
\end{lstlisting}

\subsection{Error Logging}

Use appropriate log levels for different error severities.

\begin{lstlisting}[language=C]
/* ESP_LOGE - Errors that prevent operation */
ESP_LOGE(s_TAG, "Failed to initialize bus '%s': %s",
         config->name, esp_err_to_name(ret));

/* ESP_LOGW - Warnings that don't prevent operation */
ESP_LOGW(manager->tag, "Failed to register pins for bus '%s': %s "
                       "(bus still added)", config->name, esp_err_to_name(ret));

/* ESP_LOGI - Informational messages */
ESP_LOGI(manager->tag, "Added bus config: %s (Type: %s)",
         config->name, star_bus_type_to_string(config->type));

/* ESP_LOGD - Debug messages (compiled out in release) */
ESP_LOGD(s_TAG, "Write Success: %zu bytes to '%s' (Addr 0x%02X)",
         len, config->name, address);
\end{lstlisting}

\subsection{Summary}

\begin{itemize}
    \item Return \texttt{esp\_err\_t} from all functions that can fail
    \item Use \texttt{ESP\_RETURN\_ON\_FALSE} for parameter validation
    \item Use \texttt{ESP\_GOTO\_ON\_ERROR} for cleanup patterns
    \item Always validate pointers before dereferencing
    \item Handle mutex timeouts as \texttt{ESP\_ERR\_TIMEOUT}
    \item Use dependency injection for error handlers
    \item Release resources in reverse order of acquisition
    \item Log errors with appropriate severity levels
\end{itemize}

